\documentclass[11pt]{article}

% ------------------------------------------------------------------
%  Weiss (ICCV 2001) -- "Deriving intrinsic images from image sequences"
%  One-page algorithm reference + symbol legend + the online (exam Q3) variant.
%  Compile with pdflatex (needs amsmath, amssymb, tcolorbox, booktabs, array).
% ------------------------------------------------------------------

\usepackage[a4paper,margin=2.1cm]{geometry}
\usepackage{amsmath,amssymb}
\usepackage{xcolor}
\usepackage[most]{tcolorbox}
\usepackage{booktabs}
\usepackage{array}
\usepackage{longtable}
\usepackage{enumitem}
\usepackage{titlesec}
\usepackage[hidelinks]{hyperref}

\setlist{nosep,leftmargin=1.4em}
\definecolor{teal}{RGB}{20,140,130}
\definecolor{amber}{RGB}{190,120,20}
\definecolor{warn}{RGB}{190,50,70}

\newtcolorbox{answer}{colback=green!6,colframe=green!55!black,boxrule=0.6pt,
  left=6pt,right=6pt,top=4pt,bottom=4pt,before skip=6pt,after skip=10pt}
\newtcolorbox{notebox}{colback=warn!6,colframe=warn!70!black,boxrule=0.6pt,
  left=6pt,right=6pt,top=4pt,bottom=4pt,before skip=6pt,after skip=10pt}
\newtcolorbox{qbox}{colback=black!3,colframe=black!25,boxrule=0.5pt,
  left=6pt,right=6pt,top=4pt,bottom=4pt,before skip=6pt,after skip=8pt}

\newcommand{\step}[1]{\smallskip\par\noindent\textbf{\textcolor{amber}{#1}}\par\nobreak}
\newcommand{\ans}[1]{\begin{answer}\textbf{Answer.} #1\end{answer}}

\titleformat{\section}{\large\bfseries}{\thesection.}{0.6em}{}

% column types for the algorithm table
\newcolumntype{L}[1]{>{\raggedright\arraybackslash}p{#1}}

\title{\textbf{Weiss, ICCV 2001}\\[2pt]\large Deriving Intrinsic Images from Image Sequences\\[2pt]
  \normalsize algorithm reference}
\date{}
\author{}

\begin{document}
\maketitle
\vspace{-2.4em}

% ==================================================================
\section{The model, and why it needs a prior}

An image is the product of what the surface \emph{is} and how the light \emph{fell}:
\[
  I(x,y,t) \;=\; R(x,y)\,L(x,y,t)
\]
with $R$ the \textbf{reflectance} (constant over time) and $L$ the \textbf{illumination} (changes every frame).
At one pixel, $T$ frames give $T$ equations in $T+1$ unknowns ($R$ and $L_1,\dots,L_T$), so the problem is
ill-posed no matter how many frames you collect --- and $R\equiv 1$, $L_t = I_t$ always fits. Worse, for
\emph{any} field $\alpha(x,y)$ the pair $R'=R\alpha$, $L'_t=L_t/\alpha$ reproduces every frame exactly.

\begin{notebox}\textbf{The prior.} Derivative filter outputs of natural images are \emph{sparse} --- mostly
zero, rarely huge --- which is well modelled by a Laplacian $P(x)=\frac1Z e^{-\alpha|x|}$. Weiss assumes this of
$f_n\star l$, the \emph{filtered} illumination. He does \textbf{not} assume $l$ itself is sparse; that weaker-looking
change is the whole contribution of the paper (see \S4).\end{notebox}

% ==================================================================
\section{Table 1 --- the original (offline) algorithm}

\renewcommand{\arraystretch}{1.45}
\begin{longtable}{@{}c L{5.9cm} L{8.2cm}@{}}
\toprule
\textbf{\#} & \textbf{Algorithm} & \textbf{In words} \\
\midrule
\endfirsthead
\toprule
\textbf{\#} & \textbf{Algorithm} & \textbf{In words} \\
\midrule
\endhead
0 & given $I_1,\dots,I_T$; \; $f_1=[\,{-1}\;\;1\,]$, \; $f_2=[\,{-1}\;\;1\,]^{\!\top}$
  & $T$ photos, static camera, only the lighting changes. Take two filters: horizontal and vertical
    derivative. \\
1 & $i_t = \log I_t$
  & Log every frame. The product $I=RL$ becomes the sum $i=r+l$: a constant we want plus a nuisance that
    varies. \\
2 & $o_n(x,y,t) = i_t \star f_n$
  & Filter every frame with both filters. \emph{Stop looking at brightness, look at edges} --- the prior is a
    statement about changes. \\
3 & $\hat r_n(x,y) = \operatorname*{median}\limits_{t} \, o_n(x,y,t)$
  & At each pixel, take the median down the time axis. A shadow edge crosses a given pixel in only a
    \emph{minority} of frames, so it is an outlier and the median discards it; a paint edge is present in
    \emph{every} frame and survives. \textbf{(Claim 1.)} \\
4 & $\displaystyle \hat r = g \star \Bigl( \sum_n f^{\,r}_n \star \hat r_n \Bigr)$
  & We now hold the \emph{edges} of the answer, not the answer. Integrate them back into an image by least
    squares (the two filters need not agree). With derivative filters this is a Poisson solve; $g$ is
    precomputed once. \\
5 & $\hat R = \exp(\hat r)$
  & Undo the log $\Rightarrow$ the reflectance image, shadows removed. \\
6 & $\hat L_t = I_t \,/\, \hat R$
  & Free of charge: divide each original frame by the reflectance to get that frame's illumination image. \\
\bottomrule
\end{longtable}
\renewcommand{\arraystretch}{1}

\vspace{-0.6em}
\noindent Steps 3 and 4 are equations (2) and (6) of the paper. The reversed filter is
$f^{\,r}_n(x,y)=f_n(-x,-y)$ and $g$ solves
\[
  g \star \Bigl( \sum_n f^{\,r}_n \star f_n \Bigr) \;=\; \delta .
\]
For derivative filters $\sum_n f^{\,r}_n \star f_n$ is the negative discrete Laplacian, so step 4 is exactly
$\nabla^2 \hat r = \operatorname{div}(\hat r_1, \hat r_2)$. \textbf{$g$ depends only on the chosen filters,
never on the data}, so it is computed once and reused forever --- which is exactly what makes Table 2 possible.

% ==================================================================
\section{Table 2 --- the online (streaming) algorithm}

Identical to Table 1 except for step 3. Steps 1, 2, 4, 5, 6 are already per-frame and pointwise; the plain
median is the \emph{only} thing that needs the whole sequence, so it is the only thing that changes.

\renewcommand{\arraystretch}{1.45}
\begin{longtable}{@{}c L{5.9cm} L{8.2cm}@{}}
\toprule
\textbf{\#} & \textbf{Algorithm} & \textbf{In words} \\
\midrule
\endfirsthead
\toprule
\textbf{\#} & \textbf{Algorithm} & \textbf{In words} \\
\midrule
\endhead
0 & \emph{once, in advance:} allocate $H_n[x,y,b]$, $b=0..255$; \; set $t=0$; \; precompute $g$
  & One 256-bin histogram per pixel per filter, plus a frame counter. This is all the memory there will ever
    be --- it does not grow. $g$ never depends on the data, so it is built once. \\
\midrule
  & \emph{for each arriving frame $I_t$:} & \\
1 & $i = \log I_t$ & Same as before. \\
2 & $o_n = i \star f_n$ & Same as before. \\
3a & $H_n\bigl[x,y,\,\mathrm{bin}(o_n)\bigr] \mathrel{+}= 1$; \quad $t \mathrel{+}= 1$
  & \textbf{The change.} Instead of storing the value, drop it into a bucket and add one to that bucket's
    tally. $O(1)$ work, and the sample itself is thrown away. \\
3b & $\hat r_n = $ smallest $b$ with $\sum_{b'\le b} H_n[x,y,b'] \ge t/2$
  & The running median, read straight off the histogram: walk up the buckets adding counts, stop when you have
    passed half the frames. Sort order is free --- it is the bin index. \\
4 & $\displaystyle \hat r = g \star \Bigl( \sum_n f^{\,r}_n \star \hat r_n \Bigr)$ & Same as before, using the
    \emph{current} medians. Available at any frame. \\
5 & $\hat R_t = \exp(\hat r)$ & Emit the reflectance image for this frame. \\
6 & $\hat L_t = I_t / \hat R_t$; \; \textbf{discard $I_t$} & Emit the illumination, then throw the frame away.
    No frame is ever stored. \\
\bottomrule
\end{longtable}
\renewcommand{\arraystretch}{1}

\vspace{-0.4em}
\noindent\begin{tabular}{@{}l l l@{}}
\toprule
 & \textbf{Table 1 (offline)} & \textbf{Table 2 (online)} \\
\midrule
Memory      & $O(t)$ frames --- unbounded          & $256 \times N$ counters per pixel --- constant in $t$ \\
Work/frame  & re-sort every sample                 & $O(1)$ per pixel per filter \\
Output      & only after the last frame            & $R_t$ every frame \\
\bottomrule
\end{tabular}

% ==================================================================
\section{Symbol legend}

\step{The five that cause the trouble}

\noindent\textbf{A filter} is a tiny grid of numbers you slide over an image, multiplying and adding as you go,
to produce a new image. It answers one fixed question at every pixel at once. Weiss's two filters ask
``how much did brightness change going right?'' and ``\dots going down?''

\medskip
\noindent$\boldsymbol{f_n}$ --- \textbf{the $n$-th filter itself}, i.e.\ those numbers. Not a logarithm of
anything; there is no capital $F$ in the paper, and the subscript indexes \emph{which} filter, nothing more.
Weiss uses $N=2$ throughout:
\[
  f_1 = [\,{-1}\;\;\;1\,] \quad(\text{horizontal}), \qquad
  f_2 = \begin{bmatrix} -1 \\ 1 \end{bmatrix} \quad(\text{vertical}).
\]

\medskip
\noindent$\boldsymbol{\star}$ --- \textbf{convolution}, the act of sliding. With $f_1$ it literally reduces to
``subtract each pixel from the one to its right''.

\medskip
\noindent$\boldsymbol{o_n(x,y,t)}$ --- \textbf{the output} of filter $n$ applied to the log frame $t$
($o$ for \emph{output}). One number per pixel, per filter, per frame: \emph{how much the log image changed at
that spot, in that direction, in that frame}. This is the data the algorithm actually works on --- from step 2
onward, the brightness values are never touched again. With two filters and $T$ frames you hold $2T$ such
images.

\medskip
\noindent\textbf{The median} --- fix one pixel and one filter, then look \emph{down the time axis} at that
pixel's $T$ values of $o_n$ and take the middle one. It is chosen because the Laplacian prior turns
maximum-likelihood into minimising $\sum_t |o_n - c|$, and the minimiser of a sum of absolute deviations is the
median. Its optimality condition counts how many samples sit on each side, never how far away they are ---
which is precisely why a handful of wild values (a shadow edge sweeping past in a few frames) cannot move it.

\medskip
\noindent$\boldsymbol{g}$ --- \textbf{the undo-filter}. After step 3 we hold the \emph{edges} of the answer,
$\hat r_n$, but we want the answer $\hat r$. $g$ is the filter that reverses ``take derivatives'': it satisfies
$g \star (\sum_n f^{\,r}_n \star f_n) = \delta$, so applying it cancels the filtering out. For derivative
filters it is the inverse Laplacian, i.e.\ a Poisson solve. Because it is defined purely by $f_1$ and $f_2$, it
is computed once, in advance, forever.

\step{Everything else}

\noindent\begin{tabular}{@{}l L{13.2cm}@{}}
\toprule
\textbf{Symbol} & \textbf{Meaning} \\
\midrule
$I,\;R,\;L$   & the real images: input, reflectance, illumination \\
$i,\;r,\;l$   & their logarithms. \emph{Lowercase means log for these three letters only} --- it is not a
                general rule, which is why $f_n$ trips people up \\
$r_n$         & shorthand for $r \star f_n$; the subscript is a \emph{filter index}, not a log index \\
$f^{\,r}_n$   & the \textbf{reversed} filter $f_n(-x,-y)$. The superscript $r$ means \emph{reversed}, not
                reflectance \\
$\;\hat{\cdot}\;$ & a hat marks an \emph{estimate}: $r$ is the truth, $\hat r$ is our recovered version \\
$\delta$      & the identity filter: $\delta \star u = u$ for any $u$ \\
$N,\;T$       & number of filters ($N=2$) and number of frames \\
$H_n[x,y,b]$  & Table 2 only: how many frames so far put this pixel's filter-$n$ response into bin $b$ \\
\bottomrule
\end{tabular}

% ==================================================================
\section{Why the median of the \emph{filtered} frames, not of the frames}

\begin{qbox}
Take two pixels $P$ and $Q$ with \emph{identical} true reflectance, where $P$ is sunlit in every frame and $Q$
is shaded in every frame.
\end{qbox}

\noindent\textbf{Median of the images} gives $P$ a bright value and $Q$ a dark one --- two different answers for
the same paint, and no number of extra frames repairs it, because every sample is biased the same way. This is
the ML estimate only if $l$ itself is sparse (``most pixels are unshadowed''), which is rarely true.

\medskip
\noindent\textbf{Median of the filtered images} gives them the same answer, because a permanent illumination
offset is \emph{constant} and a constant has zero derivative --- it never enters the filtered signal at all.
This is the ML estimate if $f_n \star l$ is sparse (``the illumination is mostly smooth, with rare sharp
edges''), which usually is true. The same failure kills the temporal \textbf{min} and \textbf{mean} filters:
both need each pixel to be cleanly lit at some moment.

\begin{notebox}\textbf{Claim 2 (robustness).} The Laplacian need not hold exactly. Let
$p_\epsilon = P\bigl(|f_n \star l| < \epsilon\bigr)$ measure sparseness directly. Since the median is within
$\epsilon$ whenever more than half the samples are, the binomial formula gives
\[
  P\bigl(|\hat r_n - r_n| < \epsilon\bigr) \;\ge\; \sum_{k>T/2}^{T} \binom{T}{k} p_\epsilon^{\,k}(1-p_\epsilon)^{T-k}.
\]
The paper's example ($85\%$ of filter outputs below $1\%$ of maximum, i.e.\ $p_\epsilon=0.85$) gives at least
$0.93$ from three frames and $0.97$ from five.
\emph{Note:} the summation limits are transposed in the printed paper; the upper tail is what reproduces those
numbers.\end{notebox}

% ==================================================================
\section{The online variant (exam Q3)}

\begin{qbox}\textbf{Question.} Give an \emph{online} algorithm for the reflectance image $R_t$ from a stream
$I_1,\dots,I_t$. Store at most 256 images (or equivalent). Hint: intensities lie in $[0,255]$.\end{qbox}

\step{Where it blocks}
Steps 1, 2, 4, 5 and 6 are all per-frame and pointwise, and $g$ is precomputed. \textbf{Only step 3 is
offline}: a textbook median must keep and sort every sample, so memory grows without bound. The whole question
therefore reduces to \emph{computing a running median in constant memory}.

\step{The fix}
The hint (``values in $[0,255]$'', ``store 256 images'') is naming a \textbf{256-bin histogram}: 256 slots per
pixel, which yields a median with no samples retained, because sort order comes free from the bin index. The
resulting procedure is \textbf{Table 2} above.

\ans{Per-pixel 256-bin histogram of the filtered log-frames $\to$ running temporal median $\to$ the fixed Weiss
reconstruction. Memory $256\times(\#\text{filters})$ counters per pixel, constant in $t$; work $O(1)$ per pixel
per frame; $R_t$ emitted every frame.}

\step{Details worth a line each}
\begin{itemize}
  \item The median is taken on filter \emph{responses}, not intensities. A difference of two values from a
        256-level alphabet has ${\sim}511$ possible values, so 256 bins costs at most half a bin of
        quantisation --- negligible, and bounded.
  \item Two filters means two histograms. Counting bytes: a counter needs $\lceil\log_2 T\rceil$ bits against
        a pixel's 8, so 256 bins per filter is 256 images' worth under 8-bit accounting and two filters is
        $2\times$ the stated budget; use 128 bins each, or invoke ``or equivalent memory footprint''. Showing
        the arithmetic is the point.
  \item A 256-frame ring buffer is a legal alternative (and adapts to slow scene change) but costs full frames
        rather than counters, re-sorts per pixel per frame, and caps the evidence at 256 frames.
  \item $\hat r$ is recovered only up to an additive constant, hence $\hat R$ up to one global scale --- fix it
        by convention (e.g.\ $\max \hat R = 1$). Clamp $I$ away from $0$ before the log.
\end{itemize}

% ==================================================================
\section{Limits}

\begin{itemize}
  \item \textbf{Linear camera assumed.} $I=RL$ needs a linear response; the paper blames a webcam's automatic
        gain control for its weaker outdoor results.
  \item \textbf{Independence assumed} over space and time, which is plainly false: shadow edges are long and
        connected, and consecutive frames are similar.
  \item \textbf{Static illumination is invisible.} A shadow that never moves across the sequence is
        mathematically indistinguishable from dark paint and stays baked into $\hat R$.
  \item \textbf{A sequence is required.} The single-image decomposition of Barrow and Tenenbaum remains open.
\end{itemize}

\end{document}
